\documentclass[11pt]{article}
\usepackage[T1]{fontenc}
\usepackage[utf8]{inputenc}
\usepackage{lmodern}
\usepackage{amsmath,amssymb}
\usepackage{longtable,booktabs,array}
% Compatibility with Pandoc's captionless tables on older longtable versions.
\ifcsname c@none\endcsname\else\newcounter{none}\fi
\usepackage{graphicx}
\usepackage[margin=25mm]{geometry}
\usepackage{hyperref}
\hypersetup{colorlinks=true,linkcolor=blue,urlcolor=blue}
\providecommand{\tightlist}{\setlength{\itemsep}{0pt}\setlength{\parskip}{0pt}}
\providecommand{\pandocbounded}[1]{#1}
\setlength{\emergencystretch}{3em}
\setcounter{secnumdepth}{-1}
\begin{document}
\section{Zariski's connectedness theorem and Stein
factorization}\label{zariskis-connectedness-theorem-and-stein-factorization}

\emph{Written and self-checked by GPT-6.1 Sol (OpenAI), in Codex at
Ultra, October 2026. Independent teaching material is CC0. Appendix A
adapts the Stacks project authors' freely accessible constructions and
retains GFDL-1.2-or-later; see the \href{licensing.html}{component
notice}. No independent human review or formalization is claimed.}

A proper morphism can collapse a connected curve, identify several
points, or carry a field extension in its constants. Its algebra of
functions distinguishes these phenomena. Stein factorization puts the
functions into a finite intermediate scheme; the remaining map has
geometrically connected fibers. The theorem on formal functions explains
why: a disconnected fiber would supply a compatible idempotent on every
infinitesimal neighborhood, and hence an idempotent in a completed local
ring.

We use \href{the-theorem-on-formal-functions.md}{formal functions},
\href{proper-morphisms-and-coherent-direct-images.md}{proper coherent
direct images}, and
\href{base-change-and-the-grothendieck-complex.md}{flat base change}.
Relative Spec, its universal property and affine reconstruction are
proved in the existing earlier
\href{../AG-MO/AG-MO-05.html}{\emph{Affine morphisms, relative Spec, and
finite morphisms}, Theorem 1.1 and Corollary 2.2}. The earlier
\emph{Discrete valuation rings, normal rings and Serre's criterion}
supplies normal rings; the proper birational application is proved here.
Connected means nonempty throughout. Unless stated otherwise, the base
is locally Noetherian.

\subsection{1. Connected components as
idempotents}\label{connected-components-as-idempotents}

For any scheme \(T\), a decomposition into open and closed subschemes \[
T=T_1\amalg T_2
\] defines an idempotent \(e\in\Gamma(T,\mathcal O_T)\): it is one on
\(T_1\) and zero on \(T_2\). Conversely, an idempotent determines such a
decomposition. At each point, its germ in the local ring is either zero
or one. Indeed \(e(1-e)=0\), and at least one of \(e,1-e\) is a unit in
a local ring; multiplication by that unit makes the other zero. The loci
where these germs are one and zero are respectively \(D(e)\) and
\(D(1-e)\). They are disjoint opens covering \(T\), so are also closed.

Thus a nonempty scheme is connected exactly when its global function
ring has no idempotents except zero and one. This does not assert that
its ring of functions is a field, reduced, or equal to the ground field.

\textbf{Lemma 1.1 (idempotents across a thickening).} If
\(T_0\hookrightarrow T\) is a closed immersion defined by a nilpotent
ideal, restriction induces a bijection on global idempotents.

\textbf{Proof.} The two schemes have the same underlying topological
space, since every prime contains a nilpotent ideal. An open-and-closed
decomposition of one is therefore precisely an open-and-closed
decomposition of the other. Assigning its zero and one functions gives
both the lift and its uniqueness. The idempotent/decomposition
correspondence commutes with restriction. This proves the asserted
bijection even when restriction on arbitrary global functions is not
surjective. \(\square\)

This last qualification is useful. Lifting an idempotent involves a
decomposition of the space, so it avoids the obstruction that can
prevent other sections from lifting.

Now let \((A,\mathfrak m)\) be Noetherian local, let
\(f:X\to\operatorname{Spec}A\) be proper, and put \[
X_n=X\times_A\operatorname{Spec}(A/\mathfrak m^n),\qquad n\geq1.
\] Every \(X_n\) has the same underlying space as the closed fiber
\(X_1\). A decomposition of \(X_1\) gives unique idempotents \(e_n\) on
all these schemes. Their uniqueness makes them compatible under
restriction. Formal functions in degree zero identifies the resulting
algebra with \[
\widehat{\Gamma(X,\mathcal O_X)}
\cong\varprojlim_n\Gamma(X_n,\mathcal O_{X_n}).
\tag{1}
\] The isomorphism is multiplicative, since its maps are restriction of
functions. A nontrivial fiber decomposition gives a nontrivial
idempotent in (1): its first coordinate is already neither zero nor one.

\subsection{2. Zariski's connectedness
theorem}\label{zariskis-connectedness-theorem}

\textbf{Theorem 2.1.} Let \(f:X\to S\) be proper, with \(S\) locally
Noetherian. If the unit map \[
\mathcal O_S\longrightarrow f_*\mathcal O_X
\] is an isomorphism, then every fiber is geometrically connected.

\textbf{Proof of connectedness.} First the fibers are nonempty.
Properness makes \(f(X)\) closed. If its complement contained a point,
its nonempty open complement \(U\) would have \(f^{-1}U=\varnothing\),
and hence \((f_*\mathcal O_X)|_U=0\). The assumed isomorphism would make
\(\mathcal O_U=0\), impossible at any point of \(U\).

Fix \(s\in S\) and replace the base by \(\operatorname{Spec}A\), where
\(A=\mathcal O_{S,s}\). Localization is flat. Flat base change therefore
preserves the unit isomorphism, so the function algebra of the changed
proper scheme is \(A\). Formula (1) becomes \[
\widehat A\cong\varprojlim_n\Gamma(X_n,\mathcal O_{X_n}).
\] The completion of a Noetherian local ring at its maximal ideal is
local. For example, an element outside the extended maximal ideal has an
inverse modulo that ideal; the successive corrections to this inverse
converge in the complete ring. A local ring has only the idempotents
zero and one. Lemma 1.1 would send a disconnected closed fiber to a
nontrivial idempotent in this ring. This is impossible, proving
connectedness of \(X_s\). \(\square\)

We give the extra argument for geometric connectedness rather than
treating it as automatic. A connected proper scheme over a field need
not remain connected after extension: \(\operatorname{Spec}\mathbf C\)
over \(\mathbf R\) is an immediate example.

For the proper schemes needed here, the finite-separable test has the
following direct proof. Let \(T/k\) be proper and connected and set
\(B=\Gamma(T,\mathcal O_T)\), a finite \(k\)-algebra by proper
coherence. It has no nontrivial idempotents, so its Artinian product
decomposition has one factor: \(B\) is local, and \(B_{\mathrm{red}}=L\)
is a finite field extension of \(k\). Flat base change identifies the
function algebra after every field extension \(K/k\) with
\(B\otimes_kK\). If the separable degree of \(L/k\) is larger than one,
\(L\otimes_k k^{\mathrm{sep}}\) has more than one Artinian factor: its
distinct separable embeddings give distinct maximal ideals. Its
idempotents lift uniquely across the nilpotent radical of
\(B\otimes k^{\mathrm{sep}}\), by Lemma 1.1. An idempotent uses finitely
many coefficients, so a nontrivial one already occurs over a finite
separable extension. Thus connectedness after every finite separable
extension forces \(L/k\) purely inseparable. Conversely, for purely
inseparable \(L/k\), the algebra \(L\otimes_k\overline K\) has exactly
one prime over every field extension \(K\): every element of \(L\) has a
power \(p^a\) in \(k\), and such an equation has a unique root in the
algebraic closure (in characteristic zero \(L=k\)). Hence \(L\otimes K\)
has no nontrivial idempotent, since faithful extension to
\(\overline K\) would preserve one. Nilpotent thickening by the radical
of \(B\) preserves that conclusion. Therefore \(T_K\) is connected for
every \(K\), which proves the geometric connectedness criterion in the
entire proper setting of Theorem 2.1.

\textbf{Completion of the proof of Theorem 2.1.} Let \(k'/\kappa(s)\) be
finite separable. It has a primitive element with monic irreducible
polynomial \(\overline P\). Lift its coefficients to
\(A=\mathcal O_{S,s}\), obtaining a monic polynomial \(P\), and set \[
B=A[z]/(P(z)).
\] This is a finite free \(A\)-algebra, hence flat and Noetherian. It is
local with residue field \(k'\): every maximal ideal lies over
\(\mathfrak m\), by integrality over the local ring, and
\(B/\mathfrak mB=k'\) has only one maximal ideal. Base change to
\(\operatorname{Spec}B\) preserves properness and the unit isomorphism
by flat base change. The connectedness argument just proved, applied at
its closed point, shows that \(X_s\times_{\kappa(s)}k'\) is connected.
The finite-separable-extension criterion now proves geometric
connectedness. \(\square\)

The theorem needs neither flatness of \(f\) nor reducedness of \(S\) or
\(X\). The information used is the entire algebra \(f_*\mathcal O_X\),
together with properness. No surjectivity of the ordinary comparison map
to every fiber function algebra was assumed.

\subsection{3. Constructing the Stein
factorization}\label{constructing-the-stein-factorization}

Let \(f:X\to S\) be proper. Its quasi-coherent algebra \[
\mathcal B=f_*\mathcal O_X
\] is coherent by proper coherence. Define \[
Y=\underline{\operatorname{Spec}}_S\mathcal B.
\] The evaluation map \(f^*\mathcal B\to\mathcal O_X\) gives a canonical
\(S\)-morphism \(g:X\to Y\), and we have \[
X\xrightarrow{g}Y\xrightarrow{\pi}S,\qquad f=\pi g.
\tag{2}
\] On an affine open \(U=\operatorname{Spec}A\subset S\), this is the
map to \(\operatorname{Spec}\Gamma(f^{-1}U,\mathcal O_X)\) specified by
its functions.

\textbf{Theorem 3.1 (Stein factorization).} In (2), \(\pi\) is finite,
\(g\) is proper and surjective with geometrically connected fibers, and
its unit map gives \[
g_*\mathcal O_X=\mathcal O_Y.
\] The factorization is canonical and commutes with flat base change.

\textbf{Proof.} Coherence makes \(\mathcal B\) finite as an
\(\mathcal O_S\)-module, so \(\pi\) is finite. In particular it is
separated and \(Y\) is locally Noetherian. To prove properness of \(g\),
factor it through its graph: \[
X\longrightarrow X\times_S Y\longrightarrow Y.
\] The graph is a closed immersion because \(Y\to S\) is separated; the
second map is a base change of the proper morphism \(f\). Their
composite is proper.

For affine \(U\subset S\), write \(B=\Gamma(f^{-1}U,\mathcal O_X)\). The
corresponding open in \(Y\) is \(\operatorname{Spec}B\). The
quasi-coherent sheaf \(g_*\mathcal O_X\) there has module of sections
precisely \(B\), since its inverse image is \(f^{-1}U\). Its module
structure is the canonical evaluation action of \(B\) on itself. The
affine module/sheaf correspondence therefore identifies the unit with an
isomorphism to \(\mathcal O_Y\). These identifications glue. Theorem 2.1
applied to \(g\) proves its geometrically connected, nonempty fibers,
including surjectivity.

The construction uses only the algebra \(f_*\mathcal O_X\) and its
evaluation. Under a flat change \(S_1\to S\), flat base change
identifies the new direct-image algebra with the pulled-back
\(\mathcal B\); relative Spec commutes with pullback. Evaluation also
pulls back to evaluation. Thus both maps in the factorization commute
with flat base change. \(\square\)

There is a useful uniqueness statement with a precise hypothesis.
Suppose \(f=\pi_0g_0\), where \(\pi_0:Z\to S\) is finite and the unit
\(\mathcal O_Z\to(g_0)_*\mathcal O_X\) is an isomorphism. Then \[
f_*\mathcal O_X=(\pi_0)_*\mathcal O_Z.
\] The affine reconstruction of a finite morphism gives \(Z\cong Y\),
and the maps from \(X\) agree by evaluation. Thus (2) is the unique such
factorization. Geometrically connected fibers alone are insufficient for
this uniqueness: a finite radicial field extension has geometrically
connected fibers but has more functions than the target.

\subsection{4. Reading the fibers
correctly}\label{reading-the-fibers-correctly}

\textbf{Proposition 4.1.} For \(s\in S\), the points of the finite
scheme \(Y_s\) correspond bijectively to the connected components of
\(X_s\). For a geometric point \(\overline s\), the geometric points of
\(Y_{\overline s}\) correspond to the connected components of
\(X_{\overline s}\).

\textbf{Proof.} A finite scheme over a field is the spectrum of a
finite-dimensional algebra. Such an algebra is Artinian and decomposes
into a finite product of local Artinian rings. Its points are therefore
open and closed. The base changed map \(X_s\to Y_s\) is proper and
surjective. The inverse image of each point is an open-and-closed part
of \(X_s\). Its underlying space is the fiber of \(g\) over that point
of \(Y\), with residue field unchanged; it is nonempty and connected by
Theorem 3.1. Thus these parts are exactly the connected components.
After extending the residue field to an algebraic closure, the same
argument uses geometric connectedness of the fibers of \(g\).
\(\square\)

The distinction between ordinary and geometric points matters. For
\(\operatorname{Spec}\mathbf C\to\operatorname{Spec}\mathbf R\), the
ordinary source has one component, whereas its geometric fiber has two.
The finite intermediate scheme records the field extension that produces
this splitting.

Another distinction matters even when the component count is correct.
The algebra of the finite scheme \(Y_s\) is \[
\mathcal B\otimes_{\mathcal O_S}\kappa(s),
\] and need not equal \(\Gamma(X_s,\mathcal O_{X_s})\). Theorem 3.1
guarantees flat base change, whereas taking a fiber is generally not
flat. Proposition 4.1 follows from the connected fibers of \(g\), so
remains valid without identifying these two algebras. If \(f\) is
proper, flat and of finite presentation with geometrically reduced
fibers, the stronger degree-zero base-change theorem does apply; the
preceding lesson on semicontinuity treats its connected-fiber case.

\subsection{5. The birational form of Zariski's Main
Theorem}\label{the-birational-form-of-zariskis-main-theorem}

\textbf{Theorem 5.1.} Let \(f:X\to S\) be a proper birational morphism
of integral Noetherian schemes, and suppose \(S\) is normal. Then the
unit is an isomorphism \[
\mathcal O_S\xrightarrow{\sim}f_*\mathcal O_X,
\] and all fibers of \(f\) are geometrically connected.

\textbf{Proof.} Work over a nonempty affine open
\(U=\operatorname{Spec}A\subset S\). Normality says that the domain
\(A\) is integrally closed in its fraction field \(K\). Because \(X\) is
integral and birational to \(S\), restriction to its generic point
embeds \[
B=\Gamma(f^{-1}U,\mathcal O_X)\hookrightarrow K.
\] The pullback of functions embeds \(A\) into this same field. Proper
coherence makes \(B\) a finite \(A\)-module. Every element \(b\in B\) is
integral over \(A\): multiplication by \(b\) on finitely many module
generators yields, by the determinant trick, a monic equation for \(b\).
Integral closedness therefore gives \(B\subset A\); the reverse
inclusion is the unit map. Thus \(A=B\), with equality induced by that
map. These local conclusions glue to the sheaf isomorphism. Theorem 2.1
now gives geometric connectedness. \(\square\)

Equivalently, the finite part of Stein factorization is a finite
birational morphism to a normal scheme and hence an isomorphism. On
affine opens this is exactly the argument that an integral intermediate
ring between \(A\) and \(K\) must be \(A\). It also proves the
corresponding assertion for integral birational morphisms, without a
finiteness hypothesis.

This birational connectedness theorem does not assert that \(f\) is an
isomorphism. A blow-up can have a positive-dimensional exceptional
fiber. Adding finite fibers changes the conclusion: the finiteness
theorem proved using formal functions makes \(f\) finite, and the
preceding integral-closure argument then makes it an isomorphism.

\subsection{6. Examples and the general-base
extension}\label{examples-and-the-general-base-extension}

\textbf{A node.} Over a field of characteristic different from two,
consider \[
C=\operatorname{Spec}k[x,y]/(y^2-x^2(x+1)).
\] Its normalization is \(\mathbf A^1_k\), with parameter \(u\) and \[
x=u^2-1,\qquad y=u(u^2-1).
\] Indeed \(u=y/x\) in the function field and is integral, satisfying
\(u^2=x+1\). The ring \(k[u]\) is generated as a module over the cubic's
coordinate ring by \(1,u\); it is integrally closed and has the same
function field, hence is its integral closure. This finite map is proper
and birational. The fiber over the node \((0,0)\) consists of
\(u=1,-1\), two distinct points. Its finite Stein part is the
normalization itself, and its connected part is the identity. The
failure of normality of \(C\) is exactly what permits the disconnected
fiber in this birational example.

\textbf{The plane blow-up.} For the blow-up
\(b:\operatorname{Bl}_0\mathbf A^2\to\mathbf A^2\), the preceding lesson
computed \(b_*\mathcal O=\mathcal O_{\mathbf A^2}\) and vanishing of the
positive direct images. Its Stein intermediate scheme is \(\mathbf A^2\)
itself. The exceptional fiber is \(\mathbf P^1\), while every other
fiber is one point; all are geometrically connected. Theorem 5.1
supplies the same connectedness conclusion because the plane is normal.

\textbf{Squaring.} The finite map \(\mathbf A^1_k\to\mathbf A^1_k\)
given by \(t=u^2\) has function algebra \(k[u]\) over \(k[t]\). Its
Stein factorization is the identity on the source followed by this
finite map. In characteristic different from two, a nonzero geometric
fiber consists of two points; the fiber over zero has the local algebra
\(k[u]/(u^2)\) and one point. In characteristic two every geometric
fiber has one point, with multiplicity two. The finite algebra records
both splitting and nilpotents, beyond the set of connected components.

\textbf{Theorem 6.1 (Stein factorization over an arbitrary base).} Let
\(f:X\to S\) be proper, with \(S\) any scheme. Put \[
\mathcal B=f_*\mathcal O_X,
\qquad Y=\underline{\operatorname{Spec}}_S\mathcal B.
\] Evaluation gives a canonical factorization \[
X\xrightarrow{g}Y\xrightarrow{\pi}S.
\] The map \(\pi\) is integral; \(g\) is proper and surjective with
geometrically connected fibers; and its unit is an isomorphism
\(\mathcal O_Y\xrightarrow{\sim}g_*\mathcal O_X\). Formation of this
factorization commutes with flat base change.

\textbf{Proof.} The direct-image algebra is quasi-coherent by the finite
affine-cover calculation in A.1. On every affine base open
\(U=\operatorname{Spec}A\), Lemma B.1 shows that
\(\Gamma(f^{-1}U,\mathcal O_X)\) is integral over \(A\). Relative Spec
therefore makes \(\pi\) an integral morphism. In particular it is affine
and separated.

The evaluation map \(f^*\mathcal B\to\mathcal O_X\) constructs \(g\).
Factor it through its graph \(X\to X\times_SY\to Y\). The graph is a
closed immersion because \(Y\to S\) is separated, and the second map is
the base change of the proper \(f\). Their composite is proper. This
argument uses the finite-type nature of closed immersions and of \(f\),
and does not require \(\pi\) to be of finite type.

Over \(U\), put \(B=\Gamma(f^{-1}U,\mathcal O_X)\). For each \(b\in B\),
the inverse image of the principal open
\(D(b)\subset\operatorname{Spec}B\) is the invertibility open
\((f^{-1}U)_b\). Formula (A.1) identifies its function algebra with
\(B_b\), with the evaluation module action. Since the principal opens
form a basis, the unit identifies \(g_*\mathcal O_X\) with
\(\mathcal O_Y\). Theorem C.2 now makes \(g\) surjective with
geometrically connected fibers. Finally B.3 identifies the pulled-back
direct-image algebra after any flat base change; relative Spec and
evaluation commute with arbitrary pullback by their affine algebra
constructions. Thus the factorization commutes with flat base change.
\(\square\)

The integral assertion is the correct general-base conclusion.
Noetherian coherence strengthens it to finiteness in Theorem 3.1. A
proper morphism need not be finitely presented over a non-Noetherian
base: for example, the closed immersion
\(\operatorname{Spec}(A/I)\to\operatorname{Spec}A\) is proper for every
ideal \(I\), even when \(I\) is not finitely generated. Lemma A.4.1
explicitly retains this case instead of hiding a finite-presentation
assumption in a Noetherian reduction.

The uniqueness argument in Section 3 applies to an integral intermediate
morphism as well, because an integral morphism is affine and is
reconstructed from its direct-image algebra. Its hypothesis remains the
unit isomorphism, not merely geometrically connected fibers.

\subsection{7. Exercises with solutions}\label{exercises-with-solutions}

\textbf{Exercise 7.1 (easy: a finite morphism).} Determine the Stein
factorization of a finite morphism \(f:X\to S\).

\textbf{Solution.} A finite morphism is affine, so its relative Spec
reconstruction is
\(X\cong\underline{\operatorname{Spec}}_S(f_*\mathcal O_X)\). The
evaluation map to this scheme is the identity under that isomorphism.
Thus its connected part is \(\operatorname{id}_X\), and its finite part
is \(f\). This conclusion does not require reduced fibers or flatness.

\textbf{Exercise 7.2 (easy: the node).} Explain why the normalization in
Section 6 does not contradict birational connectedness, and determine
its Stein intermediate scheme.

\textbf{Solution.} The target's coordinate ring is not integrally
closed: \(u=y/x\) satisfies \(u^2=x+1\) but is not in that ring. Thus
the normality hypothesis of Theorem 5.1 fails. The normalization is
finite, so Exercise 7.1 makes it the intermediate scheme, with identity
connected part. Its two points above the node are the two components of
that fiber.

\textbf{Exercise 7.3 (medium: when the connected part is trivial).}
Suppose \(f\) is proper over a locally Noetherian base, the unit
\(\mathcal O_S\to f_*\mathcal O_X\) is an isomorphism, and all fibers
have finitely many points. Prove that \(f\) is an isomorphism.

\textbf{Solution.} The finiteness consequence of formal functions makes
\(f\) finite. A finite morphism is affine and is reconstructed from its
direct-image algebra. That algebra, with its unit, is \(\mathcal O_S\),
whose relative Spec is \(S\). The reconstructed map is therefore the
identity of \(S\), proving the claim. Merely counting one point per
fiber would not suffice: the squaring map in characteristic two shows
why the algebra hypothesis is needed.

\textbf{Exercise 7.4 (medium: normal targets).} Let \(f:X\to Y\) be
proper birational between integral varieties over a field, with \(Y\)
normal. Prove that the fibers are connected. What stronger conclusion
holds?

\textbf{Solution.} Varieties here are integral separated schemes of
finite type over the field, so are Noetherian. On every affine open
\(\operatorname{Spec}A\subset Y\), functions upstairs form a finite
\(A\)-algebra inside \(\operatorname{Frac}A\). Normality identifies this
algebra with \(A\), as in Theorem 5.1. The unit is therefore an
isomorphism, and Theorem 2.1 proves geometric connectedness, which is
stronger than ordinary connectedness.

\textbf{Exercise 7.5 (hard: the formal-functions mechanism).} Prove
connectedness in Theorem 2.1 directly from formal functions, explaining
why no surjectivity of finite-level section maps is needed.

\textbf{Solution.} At \(s\), localize the base to
\(A=\mathcal O_{S,s}\). Flat base change gives
\(\Gamma(X_A,\mathcal O)=A\). A hypothetical nontrivial decomposition of
the closed fiber is a decomposition of every thickening's identical
underlying space. Its zero and one functions give uniquely compatible
idempotents \(e_n\). Formal functions identifies their limit with an
idempotent of \(\widehat A\). It is nontrivial because its first
coordinate is nontrivial, contradicting locality of \(\widehat A\). The
lift was made using decompositions, not by lifting arbitrary sections
through a surjective restriction map. Nonemptiness follows separately
from properness and the unit isomorphism as in Section 2.

\textbf{Exercise 7.6 (hard: inseparability and uniqueness).} In
characteristic \(p>0\), let \(L/k\) be a nontrivial finite purely
inseparable field extension. Explain why
\(\operatorname{Spec}L\to\operatorname{Spec}k\) has geometrically
connected fibers but is not its own connected part with direct-image
algebra \(k\).

\textbf{Solution.} For any algebraically closed extension \(\Omega/k\),
all elements of \(L\) have a unique image in \(\Omega\): a sufficiently
high \(p\)-power equation has a unique root. The finite algebra
\(L\otimes_k\Omega\) therefore has exactly one maximal ideal, so its
spectrum is nonempty and connected. The separable-closure criterion, or
invariance under purely inseparable extension, gives geometric
connectedness. Nevertheless its direct-image algebra is \(L\), not
\(k\). Its actual Stein factorization is the identity on
\(\operatorname{Spec}L\) followed by the finite map. Thus replacing the
unit-isomorphism condition in uniqueness by geometrically connected
fibers would give a false statement.

\subsection{Appendix A. Approximation and connectedness over an
arbitrary
base}\label{appendix-a.-approximation-and-connectedness-over-an-arbitrary-base}

The exact earlier limit inputs are
\href{../AG-GS/prerequisites/AG-MO/limits-and-noetherian-approximation.html}{\emph{Limits
of schemes and Noetherian approximation}, Theorem 1.1, Lemmas 2.1--2.3,
Theorems 4.1--4.2 and Theorem 5.1}. Its externally supplied Theorem 5.2
is not an input: the gluing proof is given in Lemma A.2.1 below. The
Chow input is
\href{../AG-GS/prerequisites/AG-MO/projective-morphisms-and-chows-lemma.html}{\emph{Projective
morphisms and Chow's lemma}, Theorem 4.1}, over an affine Noetherian
base. Proper coherence, formal functions and completion are the same
earlier proved inputs as in Sections 1--3.

\subsubsection{A.1. Two elementary facts about quasi-compact
schemes}\label{a.1.-two-elementary-facts-about-quasi-compact-schemes}

If \(T\) is quasi-compact and quasi-separated and
\(a\in\Gamma(T,\mathcal O_T)\), write \(T_a\) for its invertibility
open. Then \[
\Gamma(T,\mathcal O_T)_a=\Gamma(T_a,\mathcal O_T).
\tag{A.1}
\] Indeed, use a finite affine cover and finite affine covers of its
pairwise overlaps. Sections are the equalizer of two maps between finite
products of the chart modules. Localization commutes with finite
products and kernels, and on every chart it is ordinary ring
localization. This proves (A.1). In particular \(T_a\) is quasi-compact.

A scheme covered by finitely many affine opens \(T_{a_r}\), where all
the \(a_r\) are global functions, is canonically an open subscheme of
\(\operatorname{Spec}\Gamma(T,\mathcal O_T)\). On \(T_{a_r}\), (A.1)
identifies the canonical map with the isomorphism onto \(D(a_r)\). These
maps agree on intersections and glue to the claimed open immersion.
Notice that their union need not be the entire affine spectrum.

We also use quasi-coherence of the direct image for a quasi-compact,
quasi-separated morphism. Here is the degree-zero calculation needed in
this appendix. Over an affine target, choose a finite affine cover of
the inverse image and finite affine covers of the pairwise
intersections. Localizing the target ring localizes every term in the
equalizer for sections, and localization is exact. Thus sections on each
distinguished target open are the localization of sections on the whole
inverse image. This is precisely the affine criterion for
quasi-coherence. Kernels and images of maps of quasi-coherent modules
are quasi-coherent by the corresponding affine module calculation.

\subsubsection{A.2. Absolute approximation, with the gluing step
proved}\label{a.2.-absolute-approximation-with-the-gluing-step-proved}

\textbf{Lemma A.2.1.} Every qcqs scheme \(T\) is an inverse limit of
schemes of finite type over \(\mathbf Z\), with affine transition maps.
A prescribed finite affine cover of \(T\) can be represented eventually
by an affine cover of the stages.

\textbf{Proof.} The second assertion follows once the first has been
constructed: descend the finitely many opens by AG-MO-03 Lemma 2.1, make
each model affine by Lemma 2.3, and make them cover at a common stage by
Lemma 2.1.

We prove the first assertion by adding the members of a finite affine
cover, keeping the already added open as the inverse image of a stage
open. An affine scheme is approximated by its finitely generated
\(\mathbf Z\)-subalgebras, by AG-MO-03 Theorem 5.1. We describe the
induction step in full.

Suppose \(T=V\cup U\), where \(U=\operatorname{Spec}B\) is affine and
\(V\) already has an approximation \(V=\varprojlim V_i\). Put
\(W=U\cap V\). It is quasi-compact because \(T\) is quasi-separated. By
finite open descent, after restricting to a cofinal tail there are
quasi-compact opens \(W_i\subset V_i\), with exact inverse-image
compatibility, whose limit is \(W\).

We first arrange that the \(W_i\) are quasi-affine. A quasi-compact
quasi-affine scheme has finitely many global functions \(a_r\) whose
invertibility opens are affine and cover it: choose a principal-open
cover inside an affine scheme in which it is open. Refine so that each
\(W_{a_r}\) is contained in the inverse image of an affine chart of one
stage \(W_{i_0}\). Such a refinement is possible by intersecting the
principal opens of the ambient affine scheme with the affine chart
inverse images, then taking finitely many principal opens of that
ambient scheme inside those intersections. Their defining functions
restrict to global functions of \(W\).

By section descent, these functions come from a common stage. Eventual
containment of quasi-compact opens makes their stage invertibility opens
lie in the selected affine chart inverse images. Those inverse images
are affine, so each invertibility open is affine. Eventual equality of
opens makes them cover \(W_i\). The second fact in A.1 then identifies
\(W_i\) with an open in \(\operatorname{Spec}R_i\), where
\(R_i=\Gamma(W_i,\mathcal O_{W_i})\). Thus \(W_i\) is quasi-affine. Set
\(R=\Gamma(W,\mathcal O_W)=\varinjlim R_i\), by section descent, and
denote restriction \(B\to R\) by \(s\).

Form \[
Q_i=B\times_R R_i
=\{(b,r_i):s(b)=r_i|_W\}.
\tag{A.2}
\] Then \(\varinjlim Q_i=B\). Surjectivity follows by representing
\(s(b)\) at a stage. For injectivity, an element of \(Q_i\) with first
component zero has second component zero eventually, by the colimit
description of \(R\).

Choose \(g_1,\ldots,g_m\in B\) with
\(W=\bigcup D(g_l)\subset\operatorname{Spec}B\). Formula (A.1) gives
\(B_{g_l}=R_{s(g_l)}\). Represent each \(s(g_l)\) by compatible
functions \(g_{l,i}\in R_i\). Eventually the opens \((W_i)_{g_{l,i}}\)
are affine, cover \(W_i\), and identify with
\(D(g_{l,i})\subset\operatorname{Spec}R_i\). To check the affine
assertion, their limits are the affine opens \(D(g_l)\subset U\), so use
eventual affineness. Their cover is eventual open equality. Formula
(A.1) then identifies an affine \((W_i)_{g_{l,i}}\) with \(D(g_{l,i})\).
Put \(h_{l,i}=(g_l,g_{l,i})\in Q_i\).

The map \(Q_i\to R_i\) induces isomorphisms \[
(Q_i)_{h_{l,i}}\xrightarrow{\sim}(R_i)_{g_{l,i}}.
\tag{A.3}
\] Here is the algebra verification. For a fiber product
\(Q=B\times_R C\), let \(h=(g,c)\), and suppose \(B_g\to R_{s(g)}\) is
an isomorphism and \(c\) maps to \(s(g)\). Given \(d\in C\), its image
in \(R_{s(g)}\) can be represented by \(b/g^n\). After multiplying by a
further power of \(g\), the equality becomes an equality in \(R\). Thus
some pair \((g^q b,c^{n+q}d)\) belongs to \(Q\), and division by
\(h^{n+q}\) gives \(d\) in \(C_c\). This proves surjectivity after
localization. If \((b,d)\) maps to zero in \(C_c\), some power of \(c\)
kills \(d\); compatibility and injectivity of \(B_g\to R_{s(g)}\) then
show that some power of \(g\) kills \(b\). A common power of \(h\) kills
the pair. This proves injectivity and (A.3). Consequently \(W_i\) is an
open subscheme of \(\operatorname{Spec}Q_i\), covered there by the
\(D(h_{l,i})\).

For fixed \(i\), there is a finitely generated \(\mathbf Z\)-subalgebra
\(C\subset Q_i\) for which this same \(W_i\) is an open in
\(\operatorname{Spec}C\). Indeed each \(D(h_{l,i})\subset W_i\) is an
affine of finite type over \(\mathbf Z\), because \(W_i\) is an open in
the Noetherian finite-type scheme \(V_i\). Choose finitely many
generators of its ring \((Q_i)_{h_{l,i}}\); multiply them by
sufficiently large powers of \(h_{l,i}\) to place them in \(Q_i\). Take
\(C\) generated by these numerators and all \(h_{l,i}\). Then
\(C_{h_{l,i}}\to(Q_i)_{h_{l,i}}\) is surjective by the generators, and
injective because \(C\subset Q_i\); a relation killed after localization
in \(Q_i\) is already killed after localization in \(C\). Thus their
spectra are identical. The identifications agree on intersections, so
their unions identify \(W_i\) with \(\bigcup_lD_C(h_{l,i})\). Any larger
finitely generated subalgebra still works. The working subalgebras are
therefore a directed family with union \(Q_i\).

Glue \(V_i\) to \(\operatorname{Spec}C\) along \(W_i\); call the result
\(T_{i,C}\). It is of finite type over \(\mathbf Z\), and hence
Noetherian and qcqs. Order the pairs by \((i',C')\ge(i,C)\) when
\(i'\ge i\) and \(Q_i\to Q_{i'}\) maps \(C\) into \(C'\). This is
directed: at a common later \(i'\), the two finite sets of algebra
generators fit inside one working finite subalgebra \(C'\). Each
transition glues the given map on \(V_{i'}\) to the affine map on
\(\operatorname{Spec}C'\). Its inverse image of \(V_i\) is exactly
\(V_{i'}\), and its inverse image of \(\operatorname{Spec}C\) is exactly
\(\operatorname{Spec}C'\): both assertions follow on the affine chart
from the exact inverse-image condition for the open unions defined by
the compatible \(h_l\). The transition is affine, since these two target
opens cover and each restricted transition is affine.

The limit on the \(V_i\) side is \(V\); the limit on the affine side is
\(\operatorname{Spec}(\varinjlim C)=\operatorname{Spec}B=U\). The
overlap limit is \(W\). Scheme gluing therefore identifies the limit
with \(T=V\cup U\), and retains \(V\) as the inverse image of each
indicated stage open. This completes the induction. A directed preorder
in this construction can be replaced by its quotient under mutual
comparison; the transition maps on mutually comparable equivalent
objects are inverse isomorphisms. \(\square\)

\subsubsection{A.3. Finite-type closed ambient
schemes}\label{a.3.-finite-type-closed-ambient-schemes}

\textbf{Lemma A.3.1 (finite submodules).} On a qcqs scheme, every
quasi-coherent module is the directed union of its quasi-coherent
submodules of finite type. In particular a quasi-coherent ideal is the
directed union of finite-type quasi-coherent ideals.

\textbf{Proof.} First a finite-type quasi-coherent submodule
\(G\subset F|_W\) on a quasi-compact open \(W\subset T\) extends as a
finite-type submodule of \(F\). For an affine
\(T=\operatorname{Spec}R\), first extend it without a finiteness
condition by taking the kernel of \(F\to j_*(F|_W/G)\), where
\(j:W\hookrightarrow T\). This kernel is quasi-coherent by A.1 and
restricts to \(G\). Write it as \(\widetilde N\subset\widetilde M=F\).
Cover \(W\) by finitely many principal opens on which \(N\) localizes to
a finitely generated module. Choose finitely many numerators in \(N\)
for generators on those opens. The submodule \(N'\subset N\) they
generate restricts to \(N\) on all those opens, so
\(\widetilde{N'}|_W=G\).

For general \(T\), add finitely many affine opens to \(W\). At a step,
their intersection with the already treated quasi-compact open is
quasi-compact. Apply the affine assertion to extend the given submodule
across that affine open, then glue the two equal restrictions. Finitely
many steps give a finite-type quasi-coherent submodule of \(F\) on
\(T\). Starting with the module generated by one section on an affine
chart shows that every local section is contained in such a submodule.
Sums of two such submodules are again quasi-coherent and finite type, by
the affine calculation of images. They form a directed union equal to
\(F\). If \(F\subset\mathcal O_T\) is an ideal, its submodules are
ideals, so the final assertion follows. \(\square\)

\textbf{Lemma A.3.2 (finite-type ambient embedding).} A separated
finite-type morphism \(X\to\operatorname{Spec}A\) admits a closed
immersion into a separated scheme \(Y\) of finite presentation over
\(A\).

\textbf{Proof.} Approximate the qcqs scheme \(X\) absolutely by Lemma
A.2.1, and represent a finite affine cover \(U_a\subset X\) by an affine
cover \(U_{a,i}\subset T_i\) at the stages. Each
\(\Gamma(U_a,\mathcal O)\) is a finitely generated \(A\)-algebra. By
section descent, a finite list of its \(A\)-algebra generators comes
from \(\Gamma(U_{a,i},\mathcal O)\) at one common stage \(i\). The
canonical map \[
X\longrightarrow T_i\times_{\mathbf Z}\operatorname{Spec}A
\] has inverse image of \(U_{a,i}\times\operatorname{Spec}A\) exactly
\(U_a\), and the chart ring map
\(A\otimes_{\mathbf Z}\Gamma(U_{a,i},\mathcal O)\to\Gamma(U_a,\mathcal O)\)
is surjective. Hence it is a closed immersion. Denote the ambient scheme
by \(Y_0\); it is qcqs and of finite presentation over \(A\).

Let \(J\subset\mathcal O_{Y_0}\) define \(X\). By Lemma A.3.1 write
\(J=\bigcup J_\alpha\), with \(J_\alpha\) finite-type quasi-coherent
ideals. The schemes \(Y_\alpha=V(J_\alpha)\subset Y_0\) are of finite
presentation over \(A\), their transitions are closed immersions, and
their limit is \(X\) by the affine quotient-ring calculation.

Some \(Y_\alpha\to\operatorname{Spec}A\) is separated. To verify this
claim rather than assuming it, take a finite affine cover of a stage and
pull it back through the system. At the limit its pairwise intersections
are affine because \(X\to\operatorname{Spec}A\) is separated. They
become affine at a common stage by eventual affineness. For a pair of
chart rings at that stage, write the separatedness test as \[
R_\alpha\otimes_A R'_\alpha\longrightarrow C_\alpha
\tag{A.4}
\] for the intersection ring. At the chosen stage the affine
intersection is a quasi-compact open of a scheme of finite presentation
over \(A\), so its ring is a finite-type \(A\)-algebra. Hence (A.4) is a
finite-type algebra map as well, since its source already contains the
image of \(A\). Choose finitely many algebra generators of \(C_\alpha\)
over its source. Later intersection rings are quotients of this one,
since the transition maps are closed immersions; these same generators
therefore generate later rings over their corresponding chart-ring
tensor products. At the limit (A.4) is surjective by separatedness of
\(X\). Represent the preimages of the finite list of generators and
their equality equations at one common stage. Formula (A.4) is then
surjective there. Do this for the finitely many chart pairs. The affine
diagonal criterion proves separatedness at that stage. Choose it as
\(Y\). The closed immersion \(X\hookrightarrow Y\) remains. \(\square\)

The finite-type assertion for a quasi-compact immersion used above can
also be checked directly: a closed immersion is finite type; a
quasi-compact open immersion into an affine scheme is covered by
finitely many principal opens, whose rings are generated by one inverse.
This gives finite type locally and quasi-compactness globally.

\subsubsection{A.4. From a proper scheme to proper finitely presented
thickenings}\label{a.4.-from-a-proper-scheme-to-proper-finitely-presented-thickenings}

\textbf{Lemma A.4.1 (relative proper approximation).} If
\(X\to\operatorname{Spec}A\) is proper, then \[
X=\varprojlim_\alpha X_\alpha,
\tag{A.5}
\] where all \(X_\alpha\to\operatorname{Spec}A\) are proper and of
finite presentation, and every transition and every map
\(X\to X_\alpha\) is a closed immersion.

\textbf{Proof.} By Lemma A.3.2 choose a separated finitely presented
ambient \(Y\to\operatorname{Spec}A\) containing \(X\) closedly. First
construct a proper surjection \(q:Y'\to Y\) with an immersion
\(Y'\to\mathbf P^n_A\), with both \(Y'\) and \(q\) of finite
presentation over the indicated bases. Descend \(Y\) to a finitely
generated integer subalgebra of \(A\), by AG-MO-03 Theorems 5.1 and 4.1,
and retain separatedness using its Theorem 4.2. That base is affine
Noetherian. Apply the complete Noetherian Chow theorem AG-MO-10 Theorem
4.1 there, and pull its diagram back to \(A\). Every scheme and map in
the Noetherian diagram is of finite presentation. Properness, the
immersion, and surjectivity survive this base change. For surjectivity,
every nonempty fiber remains nonempty after field extension: on a
nonzero affine fiber ring, tensoring with a field extension is faithful.

Write \(X=V(J)\subset Y\). Lemma A.3.1 expresses \(J\) as a directed
union of finite-type ideals \(J_\alpha\). Put
\(X_\alpha=V(J_\alpha)\subset Y\), \(X'_\alpha=Y'\times_Y X_\alpha\),
and \(X'=Y'\times_Y X\). The limit of the \(X'_\alpha\) is \(X'\). Their
transitions are closed immersions and their maps to \(\mathbf P^n_A\)
are finite type. The map \(X'\to\operatorname{Spec}A\) is proper, since
\(q|_{X'}\) is proper and \(X\to\operatorname{Spec}A\) is proper. Its
immersion into \(\mathbf P^n_A\) is a closed immersion: its graph is
closed because projective space is separated, so the immersion is
proper, and a proper immersion is a closed immersion.

Consequently \(X'_\alpha\to\mathbf P^n_A\) is a closed immersion
eventually. Here is the exact limit argument. On a finite affine cover
\(V_r\) of \(\mathbf P^n_A\), the inverse-image limits are affine. Make
the stage inverse images affine by eventual affineness. Write their
chart maps \(R\to C_\alpha\). Their colimit \(R\to C\) is surjective;
transitions \(C_\alpha\to C_\beta\) are surjective because they come
from closed immersions; and \(C_\alpha\) is a finite-type \(R\)-algebra.
Choose its finitely many algebra generators. Their limit images are
elements from \(R\), so the finitely many corresponding equations hold
at one later stage. Because every later \(C_\beta\) is a quotient, those
same generator images still generate it. Thus \(R\to C_\beta\) is
surjective at that stage. A common stage handles all \(V_r\), proving
the assertion.

The resulting \(X'_\alpha\) are proper over \(A\). The maps
\(X'_\alpha\to X_\alpha\) are proper surjections. Hence
\(X_\alpha\to\operatorname{Spec}A\) is universally closed: after any
base change, the image of a closed subset \(Z\subset X_\alpha\) equals
the image of its closed inverse image in \(X'_\alpha\), which is closed.
Each \(X_\alpha\) is separated and of finite presentation over \(A\), so
universal closedness makes it proper. Restrict to the cofinal tail of
such indices. The quotient-ring limit is \(X\), and all indicated maps
are closed immersions. \(\square\)

\textbf{Lemma A.4.2 (properness at a Noetherian base stage).} Every
proper finitely presented \(Z\to\operatorname{Spec}A\) descends to a
proper morphism over a finitely generated \(\mathbf Z\)-subalgebra
\(A_i\subset A\).

\textbf{Proof.} Descend the finite presentation and then separatedness
by AG-MO-03 Theorems 4.1 and 4.2. At this Noetherian stage apply
Noetherian Chow to obtain a proper surjection \(Z'_i\to Z_i\) and an
immersion \(Z'_i\to\mathbf P^n_{A_i}\). After base change to \(A\), its
source \(Z'=Z'_i\times_{A_i}A\) is proper over \(A\), because it is
proper over the proper \(Z\). Its immersion into \(\mathbf P^n_A\) is
therefore closed by the graph argument in Lemma A.4.1. Theorem 4.2 of
AG-MO-03, applied to these finitely presented models, makes
\(Z'_j\to\mathbf P^n_{A_j}\) a closed immersion at a later stage. Hence
\(Z'_j\to\operatorname{Spec}A_j\) is proper. Its proper surjection onto
\(Z_j\), which remains separated and finitely presented, makes
\(Z_j\to\operatorname{Spec}A_j\) universally closed and therefore
proper, by the closed-subset argument in A.4.1. \(\square\)

This two-step construction is sufficient. We do not need to assert
without construction that an arbitrary finite-type proper \(X\) is the
base change of a single Noetherian model. Such an assertion would be
false without finite presentation. First take the closed-immersion limit
(A.5); only its finitely presented terms are descended to Noetherian
bases.

\subsubsection{B. General-base integrality and fiber
idempotents}\label{b.-general-base-integrality-and-fiber-idempotents}

\textbf{Lemma B.1 (functions are integral).} If
\(f:X\to\operatorname{Spec}A\) is proper, then
\(B=\Gamma(X,\mathcal O_X)\) is an integral \(A\)-algebra.

\textbf{Proof.} Let \(b\in B\). The section limit theorem applied to
(A.5) represents \(b\) by a function \(b_\alpha\) on a proper finitely
presented \(X_\alpha\) over \(A\). Descend \(X_\alpha\) to a proper
Noetherian model using Lemma A.4.2. Sections again commute with the
affine-base limit, so \(b_\alpha\) comes from a function \(b_i\) on a
proper model over some later finitely generated integer subalgebra
\(A_i\). Proper coherence makes its function algebra finite as an
\(A_i\)-module. The determinant trick gives a monic polynomial
annihilating \(b_i\): multiplication by \(b_i\) on finitely many module
generators gives a matrix of coefficients, and the adjugate identity
makes its monic characteristic polynomial kill every generator and hence
the unit. Pull that equation back to \(A\) and then to \(X\). Thus \(b\)
is integral over \(A\). \(\square\)

\textbf{Lemma B.2 (fiber idempotents lift as functions).} For a proper
morphism \(f:X\to S\), a point \(s\in S\), and an idempotent
\(e\in\Gamma(X_s,\mathcal O_{X_s})\), there is an element of the stalk
\((f_*\mathcal O_X)_s\) whose restriction to the fiber is \(e\). The
stalk element itself need not be an idempotent.

\textbf{Proof in the Noetherian case.} Replace the base by its local
ring \(A=\mathcal O_{S,s}\), using localization and the degree-zero flat
base-change calculation in B.3 below. Proper coherence makes
\(M=\Gamma(X_A,\mathcal O)\) finite over \(A\). The zero and one loci of
\(e\) partition the closed fiber into two open-and-closed pieces. Each
thickening \(X_n=X_A\times_A A/\mathfrak m^{n+1}\) has the same
underlying space, and the zero and one functions on these two pieces
give compatible idempotents \(e_n\). Formal functions identifies their
limit with an element of \(\widehat M\). Its image on the closed fiber
factors through \[
\widehat M/\mathfrak m\widehat M=M/\mathfrak mM.
\] This equality is the finite-module quotient formula of AG-CA19,
Proposition 2.2, equation (4). Choose a representative in \(M\) of that
residue class. It restricts to \(e\), proving the assertion in this
case.

\textbf{Proof over an arbitrary base.} Work on an affine neighborhood
\(\operatorname{Spec}A\) of \(s\). By (A.5),
\(X_s=\varprojlim (X_\alpha)_s\), with closed transitions. Sections
commute with the limit. Represent \(e\) at a stage and move once further
so that its equation \(e^2-e=0\) also holds. We obtain an idempotent
\(e_\alpha\) on the fiber of a proper finitely presented
\(X_\alpha\to\operatorname{Spec}A\).

Descend this morphism properly to a finitely generated integer
subalgebra \(A_i\subset A\) using A.4.2. Write \(\mathfrak p\subset A\)
for \(s\) and \(\mathfrak p_j=\mathfrak p\cap A_j\). Then \[
\kappa(s)=\varinjlim_{j\ge i}\kappa(\mathfrak p_j).
\tag{B.1}
\] Indeed the quotients \(A_j/\mathfrak p_j\) are subrings of
\(A/\mathfrak p\) with union \(A/\mathfrak p\), and every fraction uses
a numerator and a nonzero denominator lying in a common stage. The fiber
of \(X_\alpha\) is therefore the inverse limit of the Noetherian model
fibers over the \(\mathfrak p_j\), with affine transitions. Section
descent and eventual equality represent \(e_\alpha\) as an idempotent on
one such fiber. The Noetherian assertion produces a stalk element of the
model direct image mapping to that idempotent. Pullback gives a function
on a neighborhood of \(s\) in \(X_\alpha\), and restriction to \(X\)
gives a stalk element of \(f_*\mathcal O_X\). Its fiber restriction is
\(e\). \(\square\)

\textbf{Lemma B.3 (degree-zero flat base change).} For a proper morphism
\(f:X\to S\), formation of \(f_*\mathcal O_X\) commutes with flat base
change, without a Noetherian or finite-presentation assumption.

\textbf{Proof.} Over \(\operatorname{Spec}A\subset S\), its inverse
image has a finite affine cover \(U_a\). Since \(f\) is separated and
the base is affine, every \(U_a\cap U_b\) is affine: it is the inverse
image of the closed diagonal under the affine map
\(U_a\times_AU_b\to X\times_AX\). Thus \[
\Gamma(X_A,\mathcal O)
=\ker\left(\prod_a\Gamma(U_a,\mathcal O)
\longrightarrow\prod_{a,b}\Gamma(U_a\cap U_b,\mathcal O)\right).
\tag{B.2}
\] Tensoring with a flat \(A\)-algebra preserves this kernel and the
finite products. The tensor products are exactly the chart rings and
overlap rings of the base-changed cover. Hence (B.2) computes its
sections as the tensor product of the original function algebra. These
identifications localize and glue, proving the sheaf assertion.
\(\square\)

\subsubsection{C. Connectedness over arbitrary bases, including field
extensions}\label{c.-connectedness-over-arbitrary-bases-including-field-extensions}

\textbf{Lemma C.1.} A nonempty proper scheme \(T\) over a field \(k\) is
geometrically connected if it is connected after every finite separable
extension of \(k\).

\textbf{Proof.} A proper \(k\)-scheme is Noetherian, so
\(B=\Gamma(T,\mathcal O_T)\) is a nonzero finite-dimensional
\(k\)-algebra by proper coherence. An idempotent of the function algebra
is exactly a decomposition of the scheme into two open-and-closed
pieces: the local function takes the value zero or one on the respective
pieces; conversely an idempotent's invertibility open and that of its
complement form such a partition. Thus connectedness of \(T\) says that
\(B\) has no nontrivial idempotents. A finite-dimensional commutative
algebra decomposes into finitely many local algebras. Indeed every prime
quotient is a finite-dimensional domain, hence a field because
multiplication by a nonzero element is an injective and therefore
surjective linear map. There are at most \(\dim_k B\) maximal ideals:
\(r\) distinct maximal ideals give, by the Chinese remainder theorem, a
quotient with \(k\)-dimension at least \(r\). Their intersection is the
nilradical \(J\). It is generated by finitely many nilpotent elements,
since it is finite-dimensional; a sufficiently high power of \(J\) is
zero by the pigeonhole principle applied to the exponents of these
generators. The powers of the distinct maximal ideals remain pairwise
comaximal, and the intersection of common sufficiently high powers is
their product, equal to a power of \(J\), and hence zero. The Chinese
remainder theorem now decomposes \(B\) into local quotients. Since \(B\)
has no nontrivial idempotents, only one factor occurs. Consequently
\(B\) is local and \(B/J\) is a finite field extension \(L/k\).

Choose algebraic generators \(\lambda_r\) of \(L/k\). In positive
characteristic an irreducible polynomial with zero derivative is a
polynomial in \(x^p\); repeating until the derivative is nonzero shows
that some \(\lambda_r^{p^{n_r}}\) has separable minimal polynomial over
\(k\). Let \(M\subset L\) be generated by those powers. It is a finite
separable extension of \(k\), and \(L/M\) is purely inseparable. In
characteristic zero put \(M=L\). If \(M\ne k\), then
\(M\otimes_k k^{\mathrm{sep}}\) is a product of at least two copies of
\(k^{\mathrm{sep}}\): adjoin a finite list of separable generators one
at a time; in each existing factor the generator's square-free minimal
polynomial splits into distinct linear factors in the separable closure,
and the Chinese remainder theorem splits that factor accordingly.
Induction gives \([M:k]\) factors. Tensoring onward with \(L\) gives at
least two nonzero factors, since field extension is faithful. Hence
\(L\otimes_k k^{\mathrm{sep}}\) has a nontrivial idempotent. This lifts
across the nilpotent kernel of
\(B\otimes_k k^{\mathrm{sep}}\to L\otimes_k k^{\mathrm{sep}}\). One
elementary proof of nilpotent idempotent lifting is to correct \(x\)
with \(x^2-x\in I\) by Newton's formula
\(x\mapsto x-(2x-1)^{-1}(x^2-x)\): \(2x-1\) is invertible because its
square is \(1+4(x^2-x)\), and the new error lies in \(I^2\). Repeated
squaring of the error terminates for nilpotent \(I\). The lift is
nontrivial because its quotient is.

Its finitely many coefficients and its idempotent equation occur over
one finite separable extension contained in \(k^{\mathrm{sep}}\). Flat
base change B.3 would then give a nontrivial idempotent on \(T\) after
that finite separable extension, contradicting the hypothesis. Hence
\(L/k\) is purely inseparable.

For any field extension \(K/k\), pass faithfully to an algebraic closure
\(\Omega\supset K\). A finite purely inseparable extension has
generators satisfying equations \(x_r^{p^{n_r}}=a_r\in k\). Over
\(\Omega\) each polynomial has exactly one root, so all primes of
\(L\otimes_k\Omega\) contain the differences between these generators
and their unique roots. Their common quotient is \(\Omega\); the tensor
algebra is nonzero by faithfulness, so it has exactly this one prime. In
characteristic zero \(L=k\), giving the same conclusion. Thus it has no
nontrivial idempotents. Faithful scalar extension shows \(L\otimes_k K\)
has none, and lifting or reducing across the nilpotent radical shows the
same for \(B\otimes_k K\). By B.3 this is the function algebra of
\(T_K\), which is therefore connected. Nonemptiness is preserved by
field extension as above. \(\square\)

\textbf{Theorem C.2 (arbitrary-base connectedness).} If \(f:X\to S\) is
proper and the unit \(\mathcal O_S\to f_*\mathcal O_X\) is an
isomorphism, its fibers are nonempty and geometrically connected.

\textbf{Proof.} Properness makes the image closed. Its open complement,
if nonempty, would have zero direct-image structure sheaf, contradicting
the unit isomorphism at any point of that open. Thus all fibers are
nonempty.

For \(s\in S\), B.2 says that every idempotent \(e\) of the fiber
function algebra comes from \((f_*\mathcal O_X)_s=\mathcal O_{S,s}\).
The map from this local ring to the fiber function algebra factors
through \(\kappa(s)\). In a nonempty fiber the field \(\kappa(s)\)
embeds into the function algebra, since a unital map from a field to a
nonzero ring is injective. Hence an idempotent in its image is zero or
one. There are no nontrivial fiber idempotents, so \(X_s\) is connected.

Let \(k'/\kappa(s)\) be finite separable. Localize the base to
\(A=\mathcal O_{S,s}\), using B.3. Choose finitely many algebraic
generators \(\alpha_r\) of this field extension, lift the coefficients
of each generator's monic minimal polynomial to a monic
\(P_r\in A[z_r]\), and put \(C_0=A[z_1,\ldots,z_m]/(P_1,\ldots,P_m)\).
The residue classes of the monomials with each exponent smaller than
\(\deg P_r\) form an \(A\)-basis, by successive monic division; thus
\(C_0\) is finite free and flat over \(A\). Reduction modulo the maximal
ideal of \(A\) and evaluation \(z_r\mapsto\alpha_r\) give a surjection
\(C_0\to k'\). Its kernel \(N\) is maximal and lies over that maximal
ideal. The local ring \(C=(C_0)_N\) is flat over \(A\) and has residue
field \(k'\). Properness and the unit isomorphism survive base change to
\(C\), by B.3. The connectedness assertion just proved applies over this
arbitrary local ring and shows that \(X_s\times_{\kappa(s)}k'\) is
connected. Lemma C.1 completes geometric connectedness. This
construction requires no primitive-element theorem. \(\square\)

\subsection{References}\label{references}

\begin{itemize}
\tightlist
\item
  \textbf{{[}Stacks{]}} The Stacks project authors, \emph{The Stacks
  project}, in its AI Integrated Stacks Project edition: construction
  \href{https://kokunoyumeto.github.io/stacks-zh-hans-cn/en/more-morphisms.html\#more-morphisms-lemma-stein-universally-closed}{Tag
  03GY}; Noetherian Stein factorization
  \href{https://kokunoyumeto.github.io/stacks-zh-hans-cn/en/more-morphisms.html\#more-morphisms-theorem-stein-factorization-Noetherian}{Tag
  03H0}; general Stein factorization
  \href{https://kokunoyumeto.github.io/stacks-zh-hans-cn/en/more-morphisms.html\#more-morphisms-theorem-stein-factorization-general}{Tag
  03H2}.
\item
  The exact open field-extension criteria are
  \href{https://kokunoyumeto.github.io/stacks-zh-hans-cn/en/varieties.html\#varieties-lemma-characterize-geometrically-disconnected}{Tag
  0389},
  \href{https://kokunoyumeto.github.io/stacks-zh-hans-cn/en/varieties.html\#varieties-lemma-characterize-geometrically-connected}{Tag
  0387} and
  \href{https://kokunoyumeto.github.io/stacks-zh-hans-cn/en/varieties.html\#varieties-lemma-separably-closed-field-connected-components}{Tag
  0363}. Fiber idempotent lifting over arbitrary bases is
  \href{https://kokunoyumeto.github.io/stacks-zh-hans-cn/en/perfect.html\#perfect-lemma-proper-idempotent-on-fibre}{Tag
  0G7X}.
\item
  Normal targets:
  \href{https://kokunoyumeto.github.io/stacks-zh-hans-cn/en/more-morphisms.html\#more-morphisms-lemma-geometrically-connected-fibres-towards-normal}{Tag
  0AY8}, and integral birational maps
  \href{https://kokunoyumeto.github.io/stacks-zh-hans-cn/en/morphisms.html\#morphisms-lemma-finite-birational-over-normal}{Tag
  0AB1}.
\item
  Those open reference proofs retain GNU FDL 1.2. This exposition, the
  Noetherian connectedness and Stein proofs, the birational proof,
  examples and solutions are independently written CC0. The
  arbitrary-base construction and its geometric-connectedness criterion
  are proved in Appendix A and Theorem 6.1. These free references are
  construction material and supplementary reading, not substitutes for
  those arguments. The new appendix adapts the Stacks constructions and
  retains GFDL-1.2-or-later; the independently authored portions retain
  CC0.
\end{itemize}
\clearpage\section*{Copyright and licence of Stacks adaptations}
The identified Stacks proof adaptations retain the authorship of the Stacks project authors and AI Integrated Stacks Project contributors. Adaptations and added exposition by GPT-6.1 Sol (OpenAI), Codex, Ultra, 5 October 2026. Permission is granted to copy, distribute and/or modify these identified components under the GNU Free Documentation License, Version 1.2 or any later version published by the Free Software Foundation; with no Invariant Sections, no Front-Cover Texts, and no Back-Cover Texts. Independently authored portions retain their stated CC0 terms. A copy of the licence follows.

\begingroup\footnotesize
\begin{verbatim}
		GNU Free Documentation License
		  Version 1.2, November 2002


 Copyright (C) 2000,2001,2002  Free Software Foundation, Inc.
     51 Franklin St, Fifth Floor, Boston, MA  02110-1301  USA
 Everyone is permitted to copy and distribute verbatim copies
 of this license document, but changing it is not allowed.


0. PREAMBLE

The purpose of this License is to make a manual, textbook, or other
functional and useful document "free" in the sense of freedom: to
assure everyone the effective freedom to copy and redistribute it,
with or without modifying it, either commercially or noncommercially.
Secondarily, this License preserves for the author and publisher a way
to get credit for their work, while not being considered responsible
for modifications made by others.

This License is a kind of "copyleft", which means that derivative
works of the document must themselves be free in the same sense.  It
complements the GNU General Public License, which is a copyleft
license designed for free software.

We have designed this License in order to use it for manuals for free
software, because free software needs free documentation: a free
program should come with manuals providing the same freedoms that the
software does.  But this License is not limited to software manuals;
it can be used for any textual work, regardless of subject matter or
whether it is published as a printed book.  We recommend this License
principally for works whose purpose is instruction or reference.


1. APPLICABILITY AND DEFINITIONS

This License applies to any manual or other work, in any medium, that
contains a notice placed by the copyright holder saying it can be
distributed under the terms of this License.  Such a notice grants a
world-wide, royalty-free license, unlimited in duration, to use that
work under the conditions stated herein.  The "Document", below,
refers to any such manual or work.  Any member of the public is a
licensee, and is addressed as "you".  You accept the license if you
copy, modify or distribute the work in a way requiring permission
under copyright law.

A "Modified Version" of the Document means any work containing the
Document or a portion of it, either copied verbatim, or with
modifications and/or translated into another language.

A "Secondary Section" is a named appendix or a front-matter section of
the Document that deals exclusively with the relationship of the
publishers or authors of the Document to the Document's overall subject
(or to related matters) and contains nothing that could fall directly
within that overall subject.  (Thus, if the Document is in part a
textbook of mathematics, a Secondary Section may not explain any
mathematics.)  The relationship could be a matter of historical
connection with the subject or with related matters, or of legal,
commercial, philosophical, ethical or political position regarding
them.

The "Invariant Sections" are certain Secondary Sections whose titles
are designated, as being those of Invariant Sections, in the notice
that says that the Document is released under this License.  If a
section does not fit the above definition of Secondary then it is not
allowed to be designated as Invariant.  The Document may contain zero
Invariant Sections.  If the Document does not identify any Invariant
Sections then there are none.

The "Cover Texts" are certain short passages of text that are listed,
as Front-Cover Texts or Back-Cover Texts, in the notice that says that
the Document is released under this License.  A Front-Cover Text may
be at most 5 words, and a Back-Cover Text may be at most 25 words.

A "Transparent" copy of the Document means a machine-readable copy,
represented in a format whose specification is available to the
general public, that is suitable for revising the document
straightforwardly with generic text editors or (for images composed of
pixels) generic paint programs or (for drawings) some widely available
drawing editor, and that is suitable for input to text formatters or
for automatic translation to a variety of formats suitable for input
to text formatters.  A copy made in an otherwise Transparent file
format whose markup, or absence of markup, has been arranged to thwart
or discourage subsequent modification by readers is not Transparent.
An image format is not Transparent if used for any substantial amount
of text.  A copy that is not "Transparent" is called "Opaque".

Examples of suitable formats for Transparent copies include plain
ASCII without markup, Texinfo input format, LaTeX input format, SGML
or XML using a publicly available DTD, and standard-conforming simple
HTML, PostScript or PDF designed for human modification.  Examples of
transparent image formats include PNG, XCF and JPG.  Opaque formats
include proprietary formats that can be read and edited only by
proprietary word processors, SGML or XML for which the DTD and/or
processing tools are not generally available, and the
machine-generated HTML, PostScript or PDF produced by some word
processors for output purposes only.

The "Title Page" means, for a printed book, the title page itself,
plus such following pages as are needed to hold, legibly, the material
this License requires to appear in the title page.  For works in
formats which do not have any title page as such, "Title Page" means
the text near the most prominent appearance of the work's title,
preceding the beginning of the body of the text.

A section "Entitled XYZ" means a named subunit of the Document whose
title either is precisely XYZ or contains XYZ in parentheses following
text that translates XYZ in another language.  (Here XYZ stands for a
specific section name mentioned below, such as "Acknowledgements",
"Dedications", "Endorsements", or "History".)  To "Preserve the Title"
of such a section when you modify the Document means that it remains a
section "Entitled XYZ" according to this definition.

The Document may include Warranty Disclaimers next to the notice which
states that this License applies to the Document.  These Warranty
Disclaimers are considered to be included by reference in this
License, but only as regards disclaiming warranties: any other
implication that these Warranty Disclaimers may have is void and has
no effect on the meaning of this License.


2. VERBATIM COPYING

You may copy and distribute the Document in any medium, either
commercially or noncommercially, provided that this License, the
copyright notices, and the license notice saying this License applies
to the Document are reproduced in all copies, and that you add no other
conditions whatsoever to those of this License.  You may not use
technical measures to obstruct or control the reading or further
copying of the copies you make or distribute.  However, you may accept
compensation in exchange for copies.  If you distribute a large enough
number of copies you must also follow the conditions in section 3.

You may also lend copies, under the same conditions stated above, and
you may publicly display copies.


3. COPYING IN QUANTITY

If you publish printed copies (or copies in media that commonly have
printed covers) of the Document, numbering more than 100, and the
Document's license notice requires Cover Texts, you must enclose the
copies in covers that carry, clearly and legibly, all these Cover
Texts: Front-Cover Texts on the front cover, and Back-Cover Texts on
the back cover.  Both covers must also clearly and legibly identify
you as the publisher of these copies.  The front cover must present
the full title with all words of the title equally prominent and
visible.  You may add other material on the covers in addition.
Copying with changes limited to the covers, as long as they preserve
the title of the Document and satisfy these conditions, can be treated
as verbatim copying in other respects.

If the required texts for either cover are too voluminous to fit
legibly, you should put the first ones listed (as many as fit
reasonably) on the actual cover, and continue the rest onto adjacent
pages.

If you publish or distribute Opaque copies of the Document numbering
more than 100, you must either include a machine-readable Transparent
copy along with each Opaque copy, or state in or with each Opaque copy
a computer-network location from which the general network-using
public has access to download using public-standard network protocols
a complete Transparent copy of the Document, free of added material.
If you use the latter option, you must take reasonably prudent steps,
when you begin distribution of Opaque copies in quantity, to ensure
that this Transparent copy will remain thus accessible at the stated
location until at least one year after the last time you distribute an
Opaque copy (directly or through your agents or retailers) of that
edition to the public.

It is requested, but not required, that you contact the authors of the
Document well before redistributing any large number of copies, to give
them a chance to provide you with an updated version of the Document.


4. MODIFICATIONS

You may copy and distribute a Modified Version of the Document under
the conditions of sections 2 and 3 above, provided that you release
the Modified Version under precisely this License, with the Modified
Version filling the role of the Document, thus licensing distribution
and modification of the Modified Version to whoever possesses a copy
of it.  In addition, you must do these things in the Modified Version:

A. Use in the Title Page (and on the covers, if any) a title distinct
   from that of the Document, and from those of previous versions
   (which should, if there were any, be listed in the History section
   of the Document).  You may use the same title as a previous version
   if the original publisher of that version gives permission.
B. List on the Title Page, as authors, one or more persons or entities
   responsible for authorship of the modifications in the Modified
   Version, together with at least five of the principal authors of the
   Document (all of its principal authors, if it has fewer than five),
   unless they release you from this requirement.
C. State on the Title page the name of the publisher of the
   Modified Version, as the publisher.
D. Preserve all the copyright notices of the Document.
E. Add an appropriate copyright notice for your modifications
   adjacent to the other copyright notices.
F. Include, immediately after the copyright notices, a license notice
   giving the public permission to use the Modified Version under the
   terms of this License, in the form shown in the Addendum below.
G. Preserve in that license notice the full lists of Invariant Sections
   and required Cover Texts given in the Document's license notice.
H. Include an unaltered copy of this License.
I. Preserve the section Entitled "History", Preserve its Title, and add
   to it an item stating at least the title, year, new authors, and
   publisher of the Modified Version as given on the Title Page.  If
   there is no section Entitled "History" in the Document, create one
   stating the title, year, authors, and publisher of the Document as
   given on its Title Page, then add an item describing the Modified
   Version as stated in the previous sentence.
J. Preserve the network location, if any, given in the Document for
   public access to a Transparent copy of the Document, and likewise
   the network locations given in the Document for previous versions
   it was based on.  These may be placed in the "History" section.
   You may omit a network location for a work that was published at
   least four years before the Document itself, or if the original
   publisher of the version it refers to gives permission.
K. For any section Entitled "Acknowledgements" or "Dedications",
   Preserve the Title of the section, and preserve in the section all
   the substance and tone of each of the contributor acknowledgements
   and/or dedications given therein.
L. Preserve all the Invariant Sections of the Document,
   unaltered in their text and in their titles.  Section numbers
   or the equivalent are not considered part of the section titles.
M. Delete any section Entitled "Endorsements".  Such a section
   may not be included in the Modified Version.
N. Do not retitle any existing section to be Entitled "Endorsements"
   or to conflict in title with any Invariant Section.
O. Preserve any Warranty Disclaimers.

If the Modified Version includes new front-matter sections or
appendices that qualify as Secondary Sections and contain no material
copied from the Document, you may at your option designate some or all
of these sections as invariant.  To do this, add their titles to the
list of Invariant Sections in the Modified Version's license notice.
These titles must be distinct from any other section titles.

You may add a section Entitled "Endorsements", provided it contains
nothing but endorsements of your Modified Version by various
parties--for example, statements of peer review or that the text has
been approved by an organization as the authoritative definition of a
standard.

You may add a passage of up to five words as a Front-Cover Text, and a
passage of up to 25 words as a Back-Cover Text, to the end of the list
of Cover Texts in the Modified Version.  Only one passage of
Front-Cover Text and one of Back-Cover Text may be added by (or
through arrangements made by) any one entity.  If the Document already
includes a cover text for the same cover, previously added by you or
by arrangement made by the same entity you are acting on behalf of,
you may not add another; but you may replace the old one, on explicit
permission from the previous publisher that added the old one.

The author(s) and publisher(s) of the Document do not by this License
give permission to use their names for publicity for or to assert or
imply endorsement of any Modified Version.


5. COMBINING DOCUMENTS

You may combine the Document with other documents released under this
License, under the terms defined in section 4 above for modified
versions, provided that you include in the combination all of the
Invariant Sections of all of the original documents, unmodified, and
list them all as Invariant Sections of your combined work in its
license notice, and that you preserve all their Warranty Disclaimers.

The combined work need only contain one copy of this License, and
multiple identical Invariant Sections may be replaced with a single
copy.  If there are multiple Invariant Sections with the same name but
different contents, make the title of each such section unique by
adding at the end of it, in parentheses, the name of the original
author or publisher of that section if known, or else a unique number.
Make the same adjustment to the section titles in the list of
Invariant Sections in the license notice of the combined work.

In the combination, you must combine any sections Entitled "History"
in the various original documents, forming one section Entitled
"History"; likewise combine any sections Entitled "Acknowledgements",
and any sections Entitled "Dedications".  You must delete all sections
Entitled "Endorsements".


6. COLLECTIONS OF DOCUMENTS

You may make a collection consisting of the Document and other documents
released under this License, and replace the individual copies of this
License in the various documents with a single copy that is included in
the collection, provided that you follow the rules of this License for
verbatim copying of each of the documents in all other respects.

You may extract a single document from such a collection, and distribute
it individually under this License, provided you insert a copy of this
License into the extracted document, and follow this License in all
other respects regarding verbatim copying of that document.


7. AGGREGATION WITH INDEPENDENT WORKS

A compilation of the Document or its derivatives with other separate
and independent documents or works, in or on a volume of a storage or
distribution medium, is called an "aggregate" if the copyright
resulting from the compilation is not used to limit the legal rights
of the compilation's users beyond what the individual works permit.
When the Document is included in an aggregate, this License does not
apply to the other works in the aggregate which are not themselves
derivative works of the Document.

If the Cover Text requirement of section 3 is applicable to these
copies of the Document, then if the Document is less than one half of
the entire aggregate, the Document's Cover Texts may be placed on
covers that bracket the Document within the aggregate, or the
electronic equivalent of covers if the Document is in electronic form.
Otherwise they must appear on printed covers that bracket the whole
aggregate.


8. TRANSLATION

Translation is considered a kind of modification, so you may
distribute translations of the Document under the terms of section 4.
Replacing Invariant Sections with translations requires special
permission from their copyright holders, but you may include
translations of some or all Invariant Sections in addition to the
original versions of these Invariant Sections.  You may include a
translation of this License, and all the license notices in the
Document, and any Warranty Disclaimers, provided that you also include
the original English version of this License and the original versions
of those notices and disclaimers.  In case of a disagreement between
the translation and the original version of this License or a notice
or disclaimer, the original version will prevail.

If a section in the Document is Entitled "Acknowledgements",
"Dedications", or "History", the requirement (section 4) to Preserve
its Title (section 1) will typically require changing the actual
title.


9. TERMINATION

You may not copy, modify, sublicense, or distribute the Document except
as expressly provided for under this License.  Any other attempt to
copy, modify, sublicense or distribute the Document is void, and will
automatically terminate your rights under this License.  However,
parties who have received copies, or rights, from you under this
License will not have their licenses terminated so long as such
parties remain in full compliance.


10. FUTURE REVISIONS OF THIS LICENSE

The Free Software Foundation may publish new, revised versions
of the GNU Free Documentation License from time to time.  Such new
versions will be similar in spirit to the present version, but may
differ in detail to address new problems or concerns.  See
https://www.gnu.org/copyleft/.

Each version of the License is given a distinguishing version number.
If the Document specifies that a particular numbered version of this
License "or any later version" applies to it, you have the option of
following the terms and conditions either of that specified version or
of any later version that has been published (not as a draft) by the
Free Software Foundation.  If the Document does not specify a version
number of this License, you may choose any version ever published (not
as a draft) by the Free Software Foundation.


ADDENDUM: How to use this License for your documents

To use this License in a document you have written, include a copy of
the License in the document and put the following copyright and
license notices just after the title page:

    Copyright (c)  YEAR  YOUR NAME.
    Permission is granted to copy, distribute and/or modify this document
    under the terms of the GNU Free Documentation License, Version 1.2
    or any later version published by the Free Software Foundation;
    with no Invariant Sections, no Front-Cover Texts, and no Back-Cover Texts.
    A copy of the license is included in the section entitled "GNU
    Free Documentation License".

If you have Invariant Sections, Front-Cover Texts and Back-Cover Texts,
replace the "with...Texts." line with this:

    with the Invariant Sections being LIST THEIR TITLES, with the
    Front-Cover Texts being LIST, and with the Back-Cover Texts being LIST.

If you have Invariant Sections without Cover Texts, or some other
combination of the three, merge those two alternatives to suit the
situation.

If your document contains nontrivial examples of program code, we
recommend releasing these examples in parallel under your choice of
free software license, such as the GNU General Public License,
to permit their use in free software.

\end{verbatim}\endgroup
\end{document}
